% !PRINTED 437
\weberpage{437}
%<<CONT
\[
2^{\frac{n-1}{2}} \prod\limits_{1,\,\frac{n-1}{2}}^{\nu} \cos\frac{2\nu m\pi}{n} = \pm 1, \tag{19}
\]
worin einstweilen das Vorzeichen noch unbestimmt bleibt.

Wenn wir aber in (18) $\varphi$ in $\pi:2$ übergehen lassen, so erhält die rechte Seite, die sich unter der unbestimmten Form $0:0$ darstellt, den Grenzwerth $(-1)^{\frac{n-1}{2}} n$, und wir erhalten wieder die Formel, die wir im vorigen Paragraphen für eine Primzahl $n$ abgeleitet haben, allgemein für jedes ungerade $n$
\[
2^{\frac{n-1}{2}} \prod\limits_{1,\,\frac{n-1}{2}}^{\nu} \sin\frac{2\nu m\pi}{n} = \pm \sqrt{n}, \tag{20}
\]
und für $m=1$ ist das positive Zeichen zu nehmen.

Die Function $B_n$, die in Bezug auf $x$ vom $(n-1)^{\text{ten}}$ Grade ist, verschwindet, wenn $n\varphi$, aber nicht $\varphi$, ein Vielfaches von $\pi$ ist, also für
\[
x = 2\cos\frac{h\pi}{n}, \tag{21}
\]
worin $h$ eine ganze, nicht durch $n$ theilbare Zahl ist. Nehmen wir zunächst $n$ gerade an, so erhalten wir alle in (21) enthaltenen \gesperrt{verschiedenen} Werthe, wenn wir für $m$ eine beliebige zu $n$ theilerfremde Zahl setzen und
\[
h = m,\ 2m,\ 3m \ldots (n-1)m \tag{22}
\]
annehmen; denn es ist nur dann
\[
\cos\frac{h\pi}{n} = \cos\frac{h'\pi}{n},
\]
wenn entweder $h+h'$ oder $h-h'$ ein Vielfaches von $2n$ ist, was nicht eintreten kann, wenn $h$ und $h'$ beide aus der Reihe (22) genommen sind.

Unter den Werthen (21) ist einer, nämlich der dem Werth $h = nm:2$ entsprechende, gleich Null und die übrigen sind paarweise entgegengesetzt. Hiernach ergiebt sich für ein \gesperrt{gerades} $n$
\[
B_n(x) = x \prod\limits_{1,\,\frac{n}{2}-1}^{\nu} \left(x^2 - 4\cos^2\frac{\nu m\pi}{n}\right). \tag{23}
\]

Wenn $n$ ungerade ist, so ist $B_n(x)$ nur von $x^2$ abhängig; also sind je zwei der Werthe (21) entgegengesetzt gleich (für $h$ und $h+n$). Wenn man, indem man wieder $m$ zu $n$ relativ prim voraussetzt,
%CONT>>
